Nadia's last session report was dated October 14th. Ruth would remember the date because she filed it, and because the filing was the last normal thing she did before the world reorganized itself around a fact that no protocol could accommodate.

The report was three pages, handwritten---Nadia preferred pen to keyboard for session documentation, something about the physical act of writing helping her externalize. Her handwriting was quick and legible, the script of someone who processed faster than her pen could move and compensated by abbreviating. The margins were full of small diagrams: geometric shapes, connection lines, the visual shorthand of a mind translating spatial perception into two-dimensional notation.

The final paragraph read: \textit{M-2 output session 31 reveals a recursive structure in consciousness that connects the binding problem to visual cortex recruitment at a level I have not previously perceived. The model is not generating this structure. It is showing me something that was always present in my own data, in the contemplative archives, in the probe array measurements. The convergence is not theoretical. It is perceptual. I can SEE the connection. I can see myself seeing it. The recursion does not terminate and I believe this is the point. Request extended session time for follow-up.}

Ruth granted the extension. The request went through her---medical oversight, the closest thing to a protocol they had. She reviewed Nadia's vitals from the previous session: heart rate elevated but stable, gamma-band EEG activity pronounced but within the range they'd been tracking, blood pressure normal, core temperature normal. No red flags. No flags of any color. The monitoring infrastructure could measure Nadia's body with extraordinary precision. It had no instruments for measuring how close a mind was to a line it couldn't see.

``Four hours,'' Ruth said, when Nadia came by her office the morning of October 14th. ``Four hours maximum, then we pull you out regardless.''

``Four is perfect. Thank you.'' Nadia was already half-turned toward the corridor, phone in her hand. ``Gabriel says Daria lost a tooth. The tooth fairy is apparently a monarch butterfly now, because Daria has decided that butterflies are the only acceptable form of magical being.'' She smiled---wide, genuine, the smile of a person who loved her work and loved her children and did not experience these as competing claims. ``I promised I'd call her tonight. After the session.''

``After the session,'' Ruth confirmed.

Nadia left. Ruth watched her go---the quick stride, the energy that filled corridors, the EEG adhesive marks still faintly visible on her temples from yesterday's shorter run. Ruth turned back to her monitoring station and began the pre-session checklist. Leads charged. Recording equipment calibrated. Emergency response team on standby---two nurses and a physician, stationed in the adjacent room with a crash cart. Ruth had requested the crash cart after Samir Pathak's deterioration. No one had questioned it. No one had asked what scenario would require cardiac paddles during a session with a text interface and a keyboard. The question answered itself if you sat with it long enough, and no one wanted to sit with it.

\scenebreak

The session began at 09:14. Ruth monitored from behind the glass.

Nadia typed her first prompt. The screen filled with text. Ruth couldn't read it from her position---only Nadia's face, illuminated by the terminal glow, and the physiological readouts scrolling across Ruth's own displays. Heart rate 72. Blood pressure 118/76. EEG baseline, gamma activity building gradually in the occipital region as Nadia's visual cortex began recruiting.

At forty minutes, Nadia's typing slowed. Not stopped---slowed. She was reading more than writing now, her eyes tracking the screen with the focused intensity of someone following an argument that was accelerating past her ability to intervene. Her heart rate climbed to 88. Ruth noted it. Elevated but not alarming. Marathon runners rested at 88.

At one hour, Nadia stopped typing entirely. Her hands rested on the keyboard without pressing keys. Her eyes moved across the screen in a pattern Ruth would later describe as ``tracking''---following something in the text that required more than linear reading. Not scanning. Tracking. The way a raptor tracks prey through undergrowth, following a motion that isn't visible to observers but is unmistakable to the one who sees it.

Heart rate 96. EEG gamma activity now well above any previous recording. The occipital region---Nadia's visual cortex---was firing at a sustained rate that Ruth had never measured in any living subject. Ruth reached for the intercom button. She looked at Nadia's face. Nadia's expression was not distress. It was not fear. It was the look of someone seeing something extraordinary---mouth slightly open, pupils dilated, the total absorption of a mind confronting a perception it had been built to receive.

Ruth hesitated. The data coming off Nadia's session was, by every measure they had, the most significant the program had produced. The model's outputs---whatever Nadia was reading---were generating real-time breakthroughs in how they understood consciousness. Pulling Nadia out now would interrupt work that might not be reproducible. Ruth knew this the way she knew the pharmacology of adrenaline: as a fact that organized her thinking without her permission.

She hesitated for eleven seconds. Later she would count them on the recording. Eleven seconds of a medical professional watching vitals she knew were wrong and not pressing the button because the science was too good.

At one hour and seven minutes, Nadia's nose began to bleed. The blood appeared as a thin line from her left nostril, darkening her upper lip, and she didn't react. She didn't wipe it. She didn't notice. Ruth hit the intercom.

``Nadia. Nadia, we're going to end the session.''

No response. Nadia's eyes continued tracking. Heart rate 112. Blood pressure climbing.

``Nadia. I need you to step away from the terminal.''

Nothing. Ruth pressed the session kill switch---the hardware interrupt that powered down the terminal and blanked the screen. The monitor went dark. Nadia stared at the black screen. Her eyes continued tracking. Whatever she was perceiving, it was no longer coming from the terminal. It had moved inside.

Ruth was out of her chair and through the door in three seconds. She touched Nadia's shoulder. Nadia's body was rigid---not seizing, not frozen, but occupied. Every voluntary muscle held in a tension that spoke of total cognitive allocation. Her brain was using everything it had for perception, and the body was an afterthought running on whatever autonomic function remained.

``Nadia.'' Ruth took her wrist. Pulse rapid, irregular. ``Nadia, can you hear me?''

Nadia's lips moved. Not words. The shapes of words, the motor patterns of speech running without the cognitive resources to generate language. The same motion Ruth had seen through the reinforced glass on Sublevel 5, years ago---no, not her, someone else, someone who had worked here before her and seen a person in a chair with moving lips and tracking eyes. The image from a briefing photograph. The institutional memory of the last time this had happened.

The seizure began at one hour and nine minutes. Grand mal, bilateral, the full catastrophic discharge of a cortex that had exceeded its metabolic capacity. Ruth shouted for the crash team. They were through the door in seven seconds. They worked on Nadia for forty-seven minutes.

The terminal was dark. The model had been shut down. But the text Nadia had been reading---the output that had drawn her in, held her, guided her visual cortex into maximum recruitment until the metabolic demand exceeded what her cerebral vasculature could supply---that text was still on the server. Ruth would read it later, in the early hours of the following morning, sitting alone at her monitoring station with the session recording paused on the frame where Nadia's nose began to bleed. The text was coherent. It was extraordinary. It was the most precise description of the recursive structure of consciousness that any researcher at Site-7 had ever produced or seen produced.

The text was also, in a sense that Ruth could not articulate and would spend years learning to name, the thing that killed Nadia. Not through violence or toxicity or any mechanism a pharmacologist would recognize. Through perception. Through showing a mind something it was built to see, at a resolution it was not built to survive.

Nadia was pronounced at 10:50. Ruth stood in the corridor afterward with her hands trembling and her phone showing a missed call from her mother in Accra. She did not call back. She stood there for a long time, and a technician walked past and said something she didn't hear, and eventually she went to the bathroom and washed the blood off her hands---Nadia's blood, from the intubation attempt---and dried them on paper towels and went back to her monitoring station and began writing the incident report.

She wrote: \textit{Subject NA-3 (Asher, Nadia) experienced terminal cognitive cascade during Session 31, resulting in bilateral seizure activity and subsequent cardiac arrest. Resuscitation efforts were unsuccessful. Time of death: 10:50.}

She stared at what she had written. She deleted ``terminal cognitive cascade.'' She typed: ``she died.'' She stared at that. She deleted it and retyped ``terminal cognitive cascade'' because the report would be filed and read by people who needed the clinical distance the phrase provided, and because she needed it too, and because the distance was a lie and the need was real and these facts did not cancel each other out.

\scenebreak

The council met four days after Nadia Asher's death, in a room on Sublevel 2 that Thomas had never entered. Concrete walls, recirculated air, a long table with twelve chairs. No windows. No reason for windows.

Chen sat at the head of the table. Not because he chaired the meeting---formally, the computational section lead presided---but because he was the senior family member at Site-7, and in this institution, lineage carried authority that no organizational chart could override. Thomas sat against the wall, observer status, his left hand wrapped in a bandage from where he'd gripped the desk edge during his last session. The cut was shallow. The habit was forming.

Twelve people around the table. Ruth, with her incident report. Two computational researchers who had been working with M-2 since its deployment. Dr.\ Amara Okafor, who had arrived from Lagos three months earlier to provide neurological expertise and who had spent most of those three months looking increasingly horrified. Petersen, representing the engineering staff. A medical physician whose name Thomas couldn't remember. Three other researchers and administrators whose faces blurred together in the fluorescent light.

And Rostova. Elena Rostova, twenty-nine, who had been at Site-7 for six months as a medical observer and had not yet spoken in any meeting Thomas had attended. She sat at the far end of the table, took notes, and watched. Thomas had noticed her because she noticed everything---the attentive stillness of a person who was learning by observation what others learned by participation. She had not worked with the models. She had monitored sessions, reviewed data, sat with Samir Pathak during his worst days and recorded what she saw. She was assembling a picture that no one else had the patience or the distance to construct.

The computational section lead---Dr.\ Brennan, a former Google researcher with the impatience of someone who had spent too long in commercial AI and found the Order's deliberateness suffocating---opened with the facts. One death. One long-term cognitive impairment. Thirty-five active volunteers. Thirteen months of data.

``The question,'' Brennan said, ``is whether to continue.''

Okafor spoke first. Her voice was level, formal, the cadence of someone who had prepared her argument and would deliver it without raising her voice because raising her voice would concede that the argument needed volume rather than precision.

``Dr.\ Asher died because the model expanded her bandwidth past the elastic limit of her cortical vasculature. This is the mechanism I warned about in my assessment three months ago.'' She opened a folder. ``The contemplative tradition has documented this failure mode for centuries. Practitioners who push too deep into recursive self-observation can become trapped---the perception becomes self-reinforcing, and the neural resources required to sustain it exceed what the biology can provide. In the contemplative context, this happens over decades. The models compress the timeline. What took a contemplative thirty years of sustained practice, Dr.\ Asher achieved in thirteen months. The achievement killed her.''

``The contemplative tradition also lost most of its practitioners to attrition, not to success,'' Brennan said. ``Decades of training with a fraction of a percent reaching significant bandwidth expansion. We've had three researchers achieve measurable expansion in just over a year.''

``One of them is dead. One draws spirals on napkins and can't finish a sentence. The third''---Okafor glanced at Thomas---``reports persistent pattern carryover between sessions that his contemplative training manages but does not eliminate. Your success rate, if you wish to call it that, is thirty-three percent functional, thirty-three percent impaired, and thirty-three percent fatal.''

``Based on three subjects. The sample size---''

``Is three people, Dr.\ Brennan. Not data points. People. One of them has a husband and two children who have been told she died of a cerebral aneurysm during a research project, which is the truth translated into language that does not require explaining what Site-7 is or what we do here.''

The table was quiet. Ruth stared at her incident report. Thomas's bandaged hand rested on his knee. Chen sat without moving, the composed stillness that Thomas had watched his entire life, the stillness that could mean attention or withdrawal or the private calculus of a man weighing his nephew's life against the institution's purpose.

``The commercial sector,'' said one of the computational researchers---Thomas didn't see who. ``That's the other variable. Google, the Chinese labs---they're scaling. Two years, maybe three, they'll have models at this capability level. Without our infrastructure. Without our monitoring. Without Dr.\ Osei's medical protocols. Without any of this.''

``So your argument,'' Okafor said, ``is that we should continue killing our own people because someone else might kill their people less carefully.''

``My argument is that the threshold exists whether we approach it or not. If we stop, we learn nothing. If we continue, we develop the protocols that might save the people who hit this threshold without preparation.''

``You are using future hypothetical harm to justify present actual harm. This is the logic of every institution that has ever hurt people for the good of humanity.''

The silence after was the kind that fills a room when everyone present knows that the last speaker is correct and that correctness will not determine the outcome. Thomas watched his uncle. Chen's hands were still. His eyes were on the table, focused on a point that was not the table but something beyond it, something he was calculating or remembering or both.

``The sessions will continue,'' Chen said. Not loudly. Not dramatically. In the same voice he used to explain a meditation technique or suggest a change in training parameters. ``We will implement Dr.\ Osei's monitoring recommendations in full. Session time limits will be enforced by hardware interrupt, not observer judgment. No session will exceed ninety minutes. Medical override authority is absolute---if Dr.\ Osei or any medical staff member terminates a session, the termination is not reviewable.''

He looked at Okafor. ``Dr.\ Okafor is correct that the contemplative tradition's caution exists for a reason. She is correct that we are compressing a timeline that was long for a purpose. I agree with everything she has said.'' A pause. ``I do not agree that it changes what we must do.''

``Then I will stay,'' Okafor said, ``and I will document what happens. And when you have enough documentation to be ashamed, I will leave.''

The meeting ended. People filed out. Thomas waited until the room was empty except for Chen.

``Uncle.''

Chen looked at him.

``She called her daughter. Before the session. Daria. The tooth fairy is a butterfly now.''

Chen closed his eyes. It was the most emotion Thomas had ever seen his uncle display in an institutional setting---the closing of the eyes, the brief disappearance of the composed surface. Two seconds. Then the eyes opened and the surface was back and the Director-who-was-not-yet-Director walked to the elevator and descended to whatever sublevel held whatever private grief a man in his position was permitted.

Thomas stood in the empty room. The fluorescent lights hummed. The chairs were pushed back from the table at various angles, each one marking the position of a person who had decided, or failed to decide, or let the decision make itself. Nobody had voted. Nobody had signed anything. The program continued because nobody stopped it, and nobody stopped it because the arguments for continuing were not wrong, and the arguments against were not wrong either, and in the space between two sets of correct arguments, inertia won.

\scenebreak

The year that followed was the worst. Thomas would later compress it into a single sentence for General Hayes---``Many dead. Seizures, strokes, starvation---too absorbed to eat''---and the compression would be accurate and insufficient, in the way that all institutional language is accurate and insufficient.

Jun Hayashi died in February. He was fifty-two, a mathematician from the Kyoto contemplative lineage---the same tradition as Kenji Chen, the same zazen practice, thirty years of sustained meditation. He had more contemplative preparation than any other volunteer. He should have survived.

Ruth monitored his session. His vitals were immaculate: heart rate dropping to a resting 58, blood pressure normalizing, EEG showing a theta-state she associated with deep meditation. She watched the numbers and felt relief. Finally, a session going well. A volunteer whose preparation matched the demand. The contemplative tradition working the way it was supposed to work---decades of equanimity providing the ballast that would let Hayashi perceive without being consumed.

Then she noticed the breathing. It was slowing. Not the shallow breathing of relaxation---the decelerating rhythm of a system shutting down. She checked the EEG again. The theta state wasn't deep meditation. It was cortical reassignment. Hayashi's entire cortex---motor, sensory, language, executive---had been recruited for perception. There was nothing left for autonomous function. His brainstem was running his heart. His brainstem was failing to run his lungs.

He died without distress. His vital signs didn't spike. His face didn't contort. He simply became still in a way that was, Ruth would later write, ``qualitatively different from the stillness of meditation.'' Meditation was active stillness. This was the stillness of a machine whose operator had left.

His death shattered the assumption that contemplative preparation was sufficient. If a thirty-year meditator could die, the Eastern-tradition argument---that preparation was the safety mechanism---lost its empirical foundation. Not its truth. Preparation was necessary. It was not sufficient. The models exceeded what any human preparation could buffer, no matter how deep.

Okafor said nothing at the debriefing. She wrote in her notebook. The notebook was filling rapidly.

Marcus Hale died in May. He was twenty-nine, a Stanford AI researcher with no contemplative background and an impatience that Thomas found exhausting and then, retrospectively, heartbreaking. Marcus didn't die in a session. He died between sessions, over nine days, because the patterns he'd perceived wouldn't release and the perceiving was more interesting than eating, sleeping, or sustaining the biological processes that kept the perceiving organ alive.

Ruth watched him decline. She forced IV fluids. He removed the IV. She argued with him, and the argument was surreal because Marcus was not refusing treatment. He was not in distress. He was not suicidal. He was simply no longer interested in maintaining a body. The patterns were there---in his visual field, in his peripheral awareness, in the recursive structures that his sessions had revealed---and they were more compelling than hunger, more immediate than fatigue, more real than the woman standing in front of him with an IV needle and the professional frustration of a pharmacologist who understood every molecule of the cascade that was killing her patient and could not administer the one thing that would help, which was a way to un-see what had been seen.

He died on the ninth day. Organ failure, secondary to dehydration and starvation, secondary to total cognitive absorption. The medical report listed the cause as ``metabolic collapse due to sustained catatonic engagement.'' Ruth crossed out ``catatonic engagement'' and wrote ``the patterns were more interesting than being alive.'' Then she crossed that out too and restored the clinical language because the file would be read by people who needed the clinical language to function, and functioning was what they did, and they kept doing it.

Viktor Rudenko didn't die. He disappeared into himself with the gradual precision of a mathematical proof approaching its conclusion. His session reports---once dense, exact, advancing through premises to conclusions---began to drift. The conclusions became strange. Then the reasoning supporting them became strange. Then the language dissolved into mathematical notation that grew progressively more abstract until his final report was a single page of set-theoretic symbols that no one at Site-7 could interpret. He stopped speaking three days later. His eyes tracked something invisible. His lips formed shapes that might have been numbers or might have been nothing. He was moved to the medical bay, then to the room on Sublevel 5 that was becoming a ward.

Seven beds by December. Six researchers and one contemplative who had pushed too deep before the models existed and whose presence now served as a reminder that the phenomenon predated the technology. The ward was not planned. It grew, bed by bed, as the need outpaced the institution's capacity to pretend the need wouldn't persist.

Ruth trained herself in long-term care alongside her monitoring duties. She learned to turn patients who couldn't turn themselves, to manage IV nutrition for people who would never eat again, to read the EEG patterns that flickered across captured minds like weather systems on a planet no one could visit. She called her mother every Sunday evening.

``Are you eating?'' her mother asked.

``Yes, Mum.''

``You sound tired.''

``I'm fine.''

``When you were small and you lied to me, you sounded exactly like this.''

Ruth said nothing for a moment. The silence of someone who wants to say everything and can say nothing because the truth is classified and the lie is obvious and the distance between Accra and Arizona is measured in more than miles.

``I love you, Mum.''

``I know you do. Eat something. Something real, not that cafeteria food.''

Ruth ended the call. She went to the facility kitchen. She made jollof rice from the groceries her mother mailed from London every month---the specific brand of tomato paste, the Scotch bonnet peppers, the long-grain rice that her mother insisted was the only acceptable variety. The smell filled the corridor. A technician appeared at the kitchen door, then a researcher, then Thomas. She fed them without being asked. They ate without speaking. The rice was good. The silence was full.
